\documentclass[a4paper,11pt]{article}
\usepackage[utf8]{inputenc}
\usepackage[T1]{fontenc}
\usepackage[english]{babel}
\usepackage[left=2.5cm,right=2.5cm,top=2cm,bottom=2cm]{geometry}
\usepackage{hyperref}
\usepackage{xcolor}
\usepackage{microtype}
\usepackage{lmodern}

\definecolor{primary}{RGB}{0, 51, 102}
\hypersetup{colorlinks=true, urlcolor=primary}

\begin{document}
\noindent
\textbf{Loïc Aron MBASSI EWOLO} \\
ENSEA -- Double Diplôme Ingénieur \\
\href{mailto:loicaron.mbassiewolo@ensea.fr}{loicaron.mbassiewolo@ensea.fr} \quad (+33) 7 59 52 33 98 \\
\href{https://github.com/Nameless0l}{github.com/Nameless0l} \quad \href{https://mbassiloic.tech}{mbassiloic.tech}

\vspace{1em}
\noindent Cergy, le 22 mars 2026

\vspace{1em}
\noindent
\textbf{Objet : Candidature -- Développeur Cas d'Usages IA} \\
\textbf{Référence : Offre d'alternance -- Framatome}

\vspace{1.5em}
\noindent Madame, Monsieur,

\vspace{0.8em}
Transformer la stratégie digitale d'un secteur critique via des agents IA et des systèmes de recommandation robustes constitue le cœur de ma formation chez ENSEA et de ma recherche à INRIA Saclay. Vous recherchez un développeur capable de concevoir des cas d'usages IA concrets, de l'idéation à l'intégration. Mon expérience en orchestration multi-agents (Google ADK), en pipeline NLP (INRIA) et en architecture logicielle événementielle (Snappy) me permet de transformer rapidement des enjeux métier en solutions IA opérationnelles.

\vspace{0.8em}
Mon stage actuel à INRIA Saclay (équipe TRIBE, avril-août 2026) m'immerge dans un processus scientifique d'analyse de données massives et de conception de modèles de confiance interprétables. Cet environnement de recherche d'excellence m'a formé aux bonnes pratiques : rigueur expérimentale, validation statistique, documentation scientifique. En parallèle, j'ai conçu un système multi-agents agricole utilisant Google ADK et RAG (Retrieval-Augmented Generation) avec Gemini, orchestrant cinq agents IA collaboratifs pour des recommandations multi-critères. Ce projet m'a exposé aux défis réels : orchestration d'inférences, gestion du contexte, fusion de plusieurs sources.

\vspace{0.8em}
Au sein du projet Snappy (plateforme de messagerie temps réel, 4ème année ENSPY), j'ai conçu l'architecture complète d'un système de chatbots intelligents avec RAG intégré. Le système expose comment encoder la connaissance métier dans des embeddings vectoriels et générer des réponses contextuelles via LLM (Gemini Flash/Pro). Cette expérience m'a appris à naviguer les défis de productionisation : latence (cible < 100ms), qualité des réponses, gestion des erreurs. Je maîtrise le cycle complet : idéation, architecture, prototypage et intégration avec un backend réactif (Spring WebFlux, Redis, WebSocket).

\vspace{0.8em}
Je suis disponible dès septembre 2026 pour une alternance de 12 mois. Convaincu que l'IA réussie repose sur une alliance entre rigueur scientifique et pragmatisme d'ingénieur, je suis prêt à incarner cette synergie chez Framatome pour développer vos prochains cas d'usages transformateurs.

\vspace{1.5em}
\noindent Veuillez agréer, Madame, Monsieur, l'expression de mes salutations distinguées.

\vspace{2em}
\noindent \textbf{Loïc Aron MBASSI EWOLO}

\end{document}
